
\subsection*{Cross (co)products}

We here review basic facts for cross (co)products.

\subsubsection*{Ordinary cases, no braided}

Let $H$ be an ordinary biaglebra.

In the category $H-\mathrm{Mod}$, we consider an algebra object and a coalgebra object.
Note that the category of vector spaces is braided with the trivial flip operator.
Every algebra and coalgebra are bradied in this sense, though we do not add such adjective.

\begin{definition}
    A $H$-module $A$ is called a \emph{$H$-module algebra} if it admits an associative multiplication $\mu: A \otimes A \rightarrow A$ with unit $\eta: k \rightarrow A$ which are $H$-module maps:
    \[
    h \triangleright ab = (h_1 \triangleright a) (h_2 \triangleright b), \quad h \triangleright 1 = \epsilon(h).
    \]
\end{definition}

\begin{definition}
    A $H$-module $C$ is called a \emph{$H$-module coalgebra} if it admits a coassociatvive comultiplication $\Delta: C \rightarrow C \otimes C$ with counit $\epsilon: C \rightarrow k$ which are $H$-module maps:
    \[
    \Delta(h \triangleright c) = (h_1 \triangleright c_1) \otimes (h_2 \triangleright c_2), \quad \epsilon(h \triangleright c) = \epsilon(h) \epsilon(c).
    \]
\end{definition}

One intereseting fact is that
\[
C \text{ is a } H \text{-module coalgebra iff } \triangleright_C \text{ is a coalgebra map.}
\]
We don't have such statement for module algebra.
It is basically due to the fact that the coalgebra structure of $H$ get involoved into tensor product of two modules.


\subsubsection*{Braided cases}